\section*{第 \ref{ch:b1:oscillators-resonance} 章——机械振子：阻尼与共振}

\begin{solution}{exo:b1:oscillators-resonance:1}
$\omega_0 = \sqrt{80/0.20} = \qty{20}{rad/s}$，$f_0 = \qty{3.2}{Hz}$，$T_0 =
\qty{0.31}{s}$；$v_{\max} = A\omega_0 = \qty{1.0}{m/s}$；$E = \tfrac12 kA^2 =
\qty{0.10}{J}$。
\end{solution}

\begin{solution}{exo:b1:oscillators-resonance:2}
$T = 2\pi\sqrt{\ell/g}$：地球上为 \qty{2.0}{s}，月球上为 \qty{4.9}{s}。
质量–弹簧的周期 $2\pi\sqrt{m/k}$ 不含 $g$：
不变。
\end{solution}

\begin{solution}{exo:b1:oscillators-resonance:3}
$\tau = 2Q/\omega_0 = 100/(4\pi) = \qty{8.0}{s}$；能量 $\propto \eu^{-2t/\tau}$
在 $t = \tau\ln2/2 = \qty{2.8}{s}$ 时减半；能量降至 $5\%$ 约需 $3\tau = \qty{24}{s}$，
约为 $48 \approx Q$ 次振动。
\end{solution}

\begin{solution}{exo:b1:oscillators-resonance:4}
$X_0 = F_0/k = \qty{1.0}{cm}$；共振时 $\approx QX_0 = \qty{10}{cm}$；
在 $3\omega_0$ 处：$X_0/\sqrt{64 + 0.09} = \qty{0.13}{cm}$。
\end{solution}

\begin{solution}{exo:b1:oscillators-resonance:5}
$v = Lf_0 = 8.0 \times 1.3 = \qty{10}{m/s}$（\qty{37}{km/h}）。在
\qty{60}{km/h}时：$f = 16.7/8.0 = \qty{2.1}{Hz}$，$x_r = 1.6$：$X/X_g =
\sqrt{(1 + 5.2)/(2.4 + 5.2)} = 0.90$——隔振效果很差；隔振
需要 $x_r \gg 1$（在 \qty{120}{km/h} 时，$x_r = 3.2$：$0.45$）。
\end{solution}

\begin{solution}{exo:b1:oscillators-resonance:6}
推力做功为正，仅当推力沿运动方向作用时；
若每个周期都施力，则所有推力的作用相加；在其他节奏下，它们
交替正负。稳态时：$2\pi E/Q = \qty{25}{J}$，$E = \qty{80}{J}$；
$mg\ell(1 - \cos\theta_{\max}) = 80$：$1 - \cos\theta_{\max} = 0.109$，
$\theta_{\max} = 27^\circ$。
\end{solution}

\begin{solution}{exo:b1:oscillators-resonance:7}
$\delta = \ln(1/0.9) = 0.105$，$Q = \pi/\delta = 30$；$\Delta f = f_0/Q =
\qty{0.067}{Hz}$。
\end{solution}

\begin{solution}{exo:b1:oscillators-resonance:8}
位移：远低于共振频率时同相（$\varphi = 0$），共振时为 $-\pi/2$，
远高于共振频率时为 $-\pi$（反相）。共振时速度与驱动力
同相。功率为 $\vect F\cdot\vect v$：
决定驱动力做功多少的，是速度的\omterm{def:b1:wave-propagation:signal}{相位}，
而非位移的\omterm{def:b1:wave-propagation:signal}{相位}。
\end{solution}

\begin{solution}{exo:b1:oscillators-resonance:9}
$V = \omega X = \omega(F_0/m)/\allowbreak\sqrt{(\omega_0^2 - \omega^2)^2 + (\omega\omega_0/Q)^2}$；
分子分母同除以 $\omega\omega_0/Q$：$V = (F_0Q/m\omega_0)/\allowbreak
\sqrt{1 + Q^2(\omega_0/\omega - \omega/\omega_0)^2}$，且 $F_0Q/m\omega_0 = F_0/\alpha$。
平方根 $\geq 1$，等于 $1$ 恰在 $\omega = \omega_0$ 处成立。
\end{solution}

\begin{solution}{exo:b1:oscillators-resonance:10}
$P = F_0^2/2\alpha = \qty{10}{W}$；$V_{\max} = F_0/\alpha = \qty{20}{m/s}$，
$E = \tfrac12 mV_{\max}^2 = \qty{40}{J}$；$E/P = \qty{4.0}{s}$ 且 $Q/\omega_0 =
m/\alpha = \qty{4.0}{s}$。
\end{solution}

\begin{solution}{exo:b1:oscillators-resonance:11}
$3\tau = \qty{10}{s}$：$\tau = \qty{3.3}{s}$，$Q = \omega_0\tau/2 = 2\pi \times
440 \times 3.3/2 \approx 4600$；$\Delta f = f_0/Q = \qty{0.1}{Hz}$。它
几乎只响应并发出一个频率：这是一种精度达
到 $10^4$ 分之几的标准。
\end{solution}

\begin{solution}{exo:b1:oscillators-resonance:12}
$m\ddot x = mg - \rho_wgS(h + x) = -\rho_wgSx$（因为 $mg = \rho_wgSh$）：
$\omega_0^2 = \rho_wgS/m = g/h$；$T = 2\pi\sqrt{0.10/9.81} = \qty{0.63}{s}$。
质量与液体密度都只通过
平衡深度 $h$ 发挥作用。
\end{solution}

\begin{solution}{pb:b1:oscillators-resonance:1}
\textbf{1.} 受力：重力（与静态伸长所产生的弹力平衡）、弹簧力
$-ku$、阻尼器力 $-\alpha\dot u$；$m\ddot x = -ku - \alpha\dot u$，其中 $\ddot x =
\ddot x_g + \ddot u$。

\textbf{2.} $\ddot u + (\omega_0/Q)\dot u + \omega_0^2u = -\ddot x_g$，$\omega_0 =
\sqrt{k/m}$，$Q = \sqrt{km}/\alpha$。

\textbf{3.} $k = m\omega_0^2 = \pi^2 = \qty{9.9}{N/m}$；$\alpha = m\omega_0/Q =
3.14/0.70 = \qty{4.5}{N.s/m}$。

\textbf{4.} $u = -a_g/\omega_0^2$：质量块下沉一个固定量——该
仪器在零频率下即为加速度计（同时也是倾斜仪）。

\textbf{5.} 相对位移 $u(t)$。

\textbf{6.} $-\omega^2\underline U + j\omega(\omega_0/Q)\underline U + \omega_0^2\underline U
= \omega^2X_g$；两边除以 $\omega_0^2$。

\textbf{7.} $U/X_g = x_r^2/\sqrt{(1 - x_r^2)^2 + x_r^2/Q^2}$.

\textbf{8.} $U \to X_g$：质量块在\omterm{thm:b1:newton-dynamics:laws}{惯性参考系}中不动
（其惯性使它保持静止，而弹簧太软，无法拖动它）；
\omterm{def:b1:point-kinematics:frame}{参考系}绕它运动，笔记录的是地面位移
——即地震仪。

\textbf{9.} $U \approx x_r^2X_g = \omega^2X_g/\omega_0^2 = a_g/\omega_0^2$：
加速度计。

\textbf{10.} $U = QX_g = 0.7X_g$：没有峰值；两种状态平滑衔接，
仪器不会振铃。

\textbf{11.} 对 $Q = 0.7$：斜率为 $2$ 的直线上升至 $1$，
在 $x_r = 1$ 附近，随后水平于 $1$；对 $Q = 5$：形状相同，
但有高度为 $5$ 的峰，位于 $x_r = 1$ 处。

\textbf{12.} $\underline U \to -X_g$：$u = -x_g$，质量保持静止，
而\omterm{def:b1:point-kinematics:frame}{参考系}在运动。

\textbf{13.} $\Delta\ell = g/\omega_0^2 = 9.81/(2\pi \times 0.02)^2 = \qty{620}{m}$：
荒谬。

\textbf{14.} $\ell = g/\omega_0^2 = \qty{620}{m}$。

\textbf{15.} 花园门：摆绕近似竖直的
轴摆动，只有 $g$ 的极小一部分使其复位，短摆臂便给出长周期。力反馈：
线圈使质量保持固定，所需电流与地面加速度成正比——
劲度系数是电子式的，
且可以\omterm{def:b1:optical-instruments:eye}{调节}。

\textbf{16.} $\Delta\ell = 9.81/\pi^2 = \qty{0.99}{m}$：可行。

\textbf{17.} $x_r = 0.2$：$U/X_g = 0.04/\sqrt{0.92 + 0.08} = 0.040$：
\qty{40}{\micro m} —— 加速度计区域，需要放大（借助
光杠杆或电子学）。

\textbf{18.} $x_r = 10$：$100/\sqrt{9801 + 204} \approx 1.00$：\qty{1.0}{mm}，
位移记录保真。

\textbf{19.} $\omega_0 = \qty{6.3e3}{rad/s}$：$U = g/\omega_0^2 = \qty{2.5e-7}{m}$
——读作电容器间隙的变化。

\textbf{20.} $U/(a_g/\omega_0^2) = 1/\sqrt{(1 - x_r^2)^2 + x_r^2/Q^2}$；取
$Q = 0.7$ 时，它单调下降，保持在 $0.9$ 以上，直至
$x_r \approx 0.7$：约 \qty{700}{Hz}。

\textbf{21.} $\tau = 2Q/\omega_0 = \qty{0.22}{ms}$：它能无滞后地
跟随下落（历时数十毫秒）。

\textbf{22.} $\tau = 1.4/3.14 = \qty{0.45}{s}$：地面发生阶跃后，
记录在半秒内即趋于稳定，无振铃。

\textbf{23.} 若 $Q$ 较大，仪器会在每次冲击后以 $f_0$ 振铃（记录的是
仪器自身而非地面），并在 $f_0$ 附近出现峰值；
$Q \approx 0.7$ 时响应平坦且迅速。

\textbf{24.} $U = QX_g = 0.7X_g$，与理想值 $X_g$（或 $a_g/\omega_0^2
= X_g$）相比：偏低 $30\%$。

\textbf{25.} 将\omterm{def:b1:wave-propagation:signal}{信号}频率与 $f_0$ 比较：频率远高于时，
质量块为固定参考点，$u$ 即位移；远低于时，$u$
等于加速度除以 $\omega_0^2$。
\end{solution}
